% RESEARCH_DRAFT_SHA256: 832108844dc5881cc08d72d73e0222a36a33caf09efe1cd822729cb922389023
% Layout and source-reference wrappers expanded for the web renderer.
\documentclass[11pt]{article}
\usepackage[margin=25mm]{geometry}
\usepackage{fontspec,unicode-math}
\setmainfont{Latin Modern Roman}
\setmathfont{Latin Modern Math}
\usepackage{amsmath,amssymb,mathtools,xurl,hyperref}
\hypersetup{colorlinks=true,urlcolor=blue,linkcolor=blue}
\setcounter{secnumdepth}{0}
\setlength{\parindent}{0pt}
\setlength{\parskip}{5pt}
\setlength{\emergencystretch}{3em}
\newcommand{\sref}[2]{\hyperlink{source:#1:#2}{[S#2]}}
\newcommand{\cataloguescope}[1]{\paragraph{Source record.}#1}
\newcommand{\floor}[1]{\left\lfloor#1\right\rfloor}
\begin{document}
\section*{Research-attempt provenance}
This standalone research attempt retains the complete original definition body below,
including its historical source and preparation notes. Those notes describe the
earlier catalogue preparation; the new attempt and its verification are identified
separately after that body. No original ID, catalogue status, or source statement is changed.

\section{Generic Point Identification from Polynomial Measurements}\label{statement}

\subsection{Definitions and mathematical statement}

\paragraph{Shared notation.}Standard number systems are $\mathbb N,\mathbb Z,\mathbb Q,\mathbb R,\mathbb C$. Specialized notation and index ranges follow the cited source and the definitions below.


A polynomial measurement map sends an unknown point $x$ to a vector of polynomial values. Generic identification concerns points outside a specified lower-dimensional exceptional set; nongeneric measurements lie on exceptional loci of the measurement map.
\paragraph{Record scope and qualifications.}
The catalogue title is retained as metadata. Named certificates, contracts, bounds or frameworks in that title are not asserted as results of the original source extract.
The domain, measurement family, base field, and exact source meaning of generic and nongeneric are not specified.

Source-record limitation: The polynomial measurement family, domain, and generic/nongeneric exceptional loci of the cited identification theorem are absent.
\paragraph{Statement recorded in the supplied source.}
Extend the known generic-identification result from dimension two to higher dimensions. More generally, characterize polynomial measurement maps for which generic points can be identified from nongeneric measurements. \href{https://doi.org/10.4171/owr/2025/9}{[S2]}
\paragraph{Source status.}
The supplied catalogue labels this source record open; that metadata was checked against the supplied snapshot on 2026-10-11. The label is preserved for provenance and is not an independent verification of present open status or mathematical correctness. The catalogue literature assessment is dated 2026-08-22.
\subsection{Short English statement}
Extend the stated identification theorem to higher dimensions and characterize generic point recovery from nongeneric measurements.
\subsection{Sources}
\begin{itemize}
\item[S1] Ulam AI and the UnsolvedMath Contributors, UnsolvedMath: A Curated Collection of Open Mathematics Problems, catalogue entry. Catalogue curation and metadata: CC BY 4.0. The source extract is retained; definitions, English summaries and qualifications are editorial additions.
\url{https://huggingface.co/datasets/ulamai/UnsolvedMath}.
\item[S2] MFO-RIMS Tandem Workshop: Optimization, Theoretical Computer Science and Algebraic Geometry: Convexity and Beyond (2025), Oberwolfach Reports, DOI 10.4171/owr/2025/9, pp. 405--450
\url{https://doi.org/10.4171/owr/2025/9}.
\item[S3] Generic global rigidity in ell-p spaces
\url{https://arxiv.org/abs/2402.18190}.
\item[S4] Algebraic hypergraph rigidity
\url{https://arxiv.org/abs/2305.18990}.
\item Source-specific Oberwolfach Reports licences, including CC BY-SA 4.0 where applicable, are retained; see THIRD\_PARTY\_NOTICES.md.
\end{itemize}
\paragraph{Preparation and review.}
The original source and UnsolvedMath contributors retain attribution for the source extract. Codex (GPT-6; exact serving revision not exposed) assisted with definitions, English summaries, scope notes and formatting. No new human mathematical authorship or independently verified solution is asserted. Catalogue identity checked against the complete supplied registry and current repository. Distinct upstream records, including repeated questions and broad or damaged extracts, are retained for complete coverage of the source CS category. A repeated record is not a new independent conjecture. Complete mathematical-equivalence and current-literature review has not been performed. Preparation and final review follow contribution prompts C and D, with source gaps disclosed. Structural and rendering checks do not establish a new theorem or open status.
\label{end:statement}

\section{Literature and context check}
The primary report, printed pp. 428--429, places the dimension-two result in generic global rigidity for even-power $\ell_q$ distances, with a broader polynomial-hypergraph measurement framework. The abbreviated record omits that structure. The Euclidean anchored reconstruction below is an auxiliary model with known anchors; it neither proves higher-dimensional $\ell_q$ rigidity nor identifies arbitrary unanchored frameworks.
\section{Research attempt: Anchor reconstruction and a local-to-global identifiability obstruction}

Let $a_0,\ldots,a_d\in\mathbb R^d$ be known affinely independent anchors. An unknown $x$ has measurements $s_i=\|x-a_i\|_2^2$. Subtracting the equation for $i=0$ gives
\[
2(a_i-a_0)^{\mathsf T}x=\|a_i\|_2^2-\|a_0\|_2^2-(s_i-s_0),\qquad 1\leq i\leq d.
\]
The coefficient matrix is invertible by affine independence, so these equations give the unique candidate. Substituting it in all original equations checks consistency. Thus a finite linear solve exactly recovers any compatible $x$ in this anchored model, and also shows how reconstruction becomes ill-conditioned when the anchor-difference matrix approaches singularity.

Now let the affine hull of all anchors have dimension $k<d$. Translate and rotate coordinates so it is $\mathbb R^k\times\{0\}$ and write $x=(u,v)$. Every measurement equals
\[
\|u-a_i\|_2^2+\|v\|_2^2.
\]
If the anchors affinely span that $k$-space, their differences determine $u$; one remaining distance determines only $\|v\|$. For $\|v\|>0$, the compatible fiber is exactly a sphere in $\mathbb R^{d-k}$ of that radius. When $k=d-1$ there are two points, related by reflection across the anchor hyperplane; when $k\leq d-2$ there are infinitely many. Zero radius is the exceptional singleton boundary.

In particular, with exactly $d$ affinely independent anchors in a hyperplane and an $x$ outside it, the measurement map to their $d$ squared distances has full-rank Jacobian despite a two-point fiber. Its Jacobian rows are $2(x-a_i)^{\mathsf T}$. Their differences span the anchor hyperplane directions, and any one row has a nonzero perpendicular component, so together they span $\mathbb R^d$. Thus local Jacobian nonsingularity does not establish global identification.

For the source's even $q>2$ distances, subtraction of $q$th-power distances leaves higher-degree terms rather than a linear system. The preceding reconstruction cannot simply be transferred. Moreover the source identifies frameworks only up to an appropriate stabilizer, whereas the anchored model fixes an external coordinate frame. Both changes matter to any proposed genericity theorem.
\section{Verification and adversarial audit}
The linear inversion requires $d+1$ affinely independent known anchors; the two-point example uses only $d$. The complete fiber description includes the zero-radius exception. Jacobian rank was checked explicitly in the two-point case, exposing the exact unsupported local-to-global inference.
\paragraph{Audit corrections.} Distinguished local recovery, global recovery, and recovery modulo stabilizers, and did not treat Euclidean distance subtraction as valid for all even powers.
\paragraph{External computation.} None; symbolic arguments and boundary checks only.
\section{Scope and remaining gaps}
No higher-dimensional generic global $\ell_q$ rigidity characterization or general polynomial measurement criterion is proved. Missing map/domain conventions in the abbreviated record are documented using the primary context, not replaced by the easier anchored Euclidean problem.
\paragraph{Integration review.} Codex (GPT-6; exact serving revision not exposed) reviewed the written arguments and expanded layout and source-reference wrappers. Repository metadata and the production web renderer were checked using Python 3.13 and JavaScript with MathJax 3.2.2. This model-assisted review is not independent human proof verification; no PDF compilation or formal proof checking is asserted.

\section{Completion Estimate}
\noindent\textbf{COMPLETION: 0\%}
The calculations below the preserved source statement investigate a concrete route,
but establish no substantiated new progress toward the entire original target.
Known special cases and failed sufficient conditions are not credited as a solution.
\section{References}
The original source references and catalogue attribution are preserved in the
Sources section above. The following URLs were actually consulted for this attempt;
retention of an additional original reference does not claim fresh consultation of it.
\begin{itemize}
\item \url{https://ems.press/content/serial-article-files/51353} (consulted 2026-10-11). 
\end{itemize}
\section{Model and contributor disclosure}
Codex GPT-6 completed written-proof review of this attempt at its declared scope; this is model-assisted review, not independent human proof checking.
OpenAI gpt-6-astra, requested reasoning configuration ultra, generated this research
attempt on 2026-10-11. The exact serving revision is not exposed. Codex GPT-6
performed orchestration and repository-format review; its exact serving revision
is also not exposed. The mathematical route, boundary checks and final adversarial
self-audit followed contribution prompts A then B in
\texttt{docs/contribution\_prompts.md}, after reading \texttt{CONTRIBUTING.md}.
Python 3.14.2 assembled the standalone documents; any mathematical computation is
specified in Verification. No theorem prover was used. No human mathematical
authorship, substantial human proof correction, or independent human proof checking
is asserted. Model-assisted review, source retrieval, and structural checks do not
establish correctness or novelty. Full private reasoning transcripts are unavailable
as reproducibility artifacts; the complete arguments, inputs and outcomes actually
used are stated here or in the companion batch audit.
\section{Final status}
Unresolved as a new research contribution. This is a concrete route investigation,
not a claim of new complete resolution or a current-status certification. A known
answer reproduced for context does not become a new contribution. The catalogue's
source-status label is preserved; this attempt status does not assert that a
source-labelled solved question is open.
\end{document}
